\documentclass[UTF8,a4paper,11pt]{ctexart}
\usepackage{amsmath,amssymb}
\usepackage[margin=2.6cm]{geometry}
\usepackage[hidelinks]{hyperref}
\title{2011年安徽高考文理科共同导数题}
\author{三千弱水数学题库}
\date{2026年10月24日}
\begin{document}
\maketitle
\noindent\textbf{出处：}文科第18题 \url{https://github.com/DxAThing/Gaokao-Math-Problems-Compilation/blob/f1ae169b58a75b3db6e35c670b3f0d5c78718188/content/2011/anhui_liberal.tex#L447}

\noindent 理科第16题 \url{https://github.com/DxAThing/Gaokao-Math-Problems-Compilation/blob/f1ae169b58a75b3db6e35c670b3f0d5c78718188/content/2011/anhui_science.tex#L362}

\paragraph{定义域和导数}
对 $a>0$，考虑函数
\[
f(x)=\frac{e^x}{1+ax^2},\qquad x\in\mathbb R.
\]
分母对全体实数严格为正，故
\[
f'(x)=\frac{e^x\bigl[a(x-1)^2+1-a\bigr]}{(1+ax^2)^2}.
\]

\paragraph{当 $a=4/3$ 时}
此时
\[
a(x-1)^2+1-a
=\frac43\left(x-\frac12\right)\left(x-\frac32\right).
\]
导数符号依次为正、负、正，
故极大值点为 $x=1/2$，极小值点为 $x=3/2$。
对应的极值分别是
\[
f(1/2)=\frac{3\sqrt e}{4},\qquad
f(3/2)=\frac{e^{3/2}}4.
\]

\paragraph{整条实数轴上的单调性}
当 $0<a<1$ 时，$a(x-1)^2+1-a>0$；
当 $a=1$ 时，该式为 $(x-1)^2$，
只在一个点为零，函数仍严格递增。
当 $a>1$ 时，导数在 $x=0$ 处为正、在 $x=1$ 处为负，
故不可能全局单调。
于是所求为
\[
\boxed{0<a\le1}.
\]

\bigskip
\noindent\textbf{质检说明：}数学推导独立验证通过；
本源码尚未实际编译，不得标注为编译通过。
\end{document}
